%% ============================================================
%%  Results-and-Discussion.tex  — corrected image placements
%%  Images in: images/  folder
%%  Macro: \inlinefig[width]{path}{caption}{label}
%% ============================================================

\section{Results and Discussion}
\label{sec:results}

%%------------------------------------------------------------
\subsection{Regime Detection}

For U1, $I_t$ spans 0.0--1.0 (mean 0.3594, std 0.2357), with 38 Calm,
11 Elevated, and 2 Crisis windows. The Crisis rate of 3.9\% reflects
deliberate calibration: the graph penalty activates rarely and with
purpose, not persistently~\citep{ang2002,ang2004}. COVID windows
(W37--W39) register as Elevated rather than Crisis — consistent with
the rapid V-shaped recovery in technology equities compared with broader
indices~\citep{longin2001,guidolin2007}. The bimodal $I_t$ distribution
validates that $\tau_{\text{crisis}} = 0.85$ separates genuinely distinct
market states~\citep{ang2002}.

\inlinefig[0.96]{images/plot1_It_timeline_U1}
  {Composite Instability Index $I_t$ across all 51 rolling windows
   (Universe U1, 2005--2025). Shaded bands mark the three crisis
   periods: 2008 GFC (W08--W11), 2020 COVID (W37--W39), and 2022
   Rate Hike (W42--W44). W09 attains the maximum $I_t = 1.000$.
   The sigmoid threshold $\tau_{\text{crisis}} = 0.85$ (dashed line)
   classifies exactly 2 windows as Crisis~\citep{ang2002,hamilton1989}.}
  {fig:it_timeline}

\inlinefig[0.82]{images/plot2_It_distribution_U1}
  {Kernel density estimate of $I_t$ for Universe U1. The bimodal
   structure with a dominant calm-regime peak near 0.20 and a
   crisis-regime tail above 0.85 is consistent with the two-regime
   equity market structure documented by \citet{ang2002}.}
  {fig:it_distribution}

%%------------------------------------------------------------
\subsection{Contagion Graph Analysis}

Across all 51 U1 windows, Agent~2 achieves 100\% eigenvector centrality
convergence with 0\% fallback to uniform centrality~\citep{bonacich1987}.
LRCX (0.2577) and INTU (0.2553) rank highest, reflecting concentrated
Vanguard and BlackRock co-ownership, while ORCL (0.2015) ranks lowest,
consistent with its lower institutional co-ownership relative to
semiconductor and software stocks~\citep{azar2018,bebchuk2017}. The
centrality standard deviation of 0.0118 across tickers in W01 confirms
heterogeneity in the penalty signal. Vanguard holds 9.3--10.7\% of 16 of
17 U1 tickers — a level of co-ownership consistent with the anticompetitive
concentration patterns documented by \citet{azar2018}.

\inlinefig[0.96]{images/plot3_bipartite_U1}
  {Bipartite institutional ownership graph for Universe U1, Window W09
   ($I_t = 1.000$). Stock nodes (circles) are coloured by eigenvector
   centrality; institution nodes (squares) are sized by total holdings.
   Edge weight equals percentage held. Vanguard Group Inc and BlackRock
   Inc are the dominant connectors, creating the co-ownership channels
   that the G-CVaR penalty is designed to
   attenuate~\citep{azar2018,billio2012}.}
  {fig:bipartite}

\inlinefig[0.82]{images/plot4_centrality_boxplot_U1}
  {Eigenvector centrality distribution across all 51 windows for the
   17 retained U1 tickers. LRCX and INTU exhibit consistently high
   centrality; ORCL is the consistent outlier. Box width encodes the
   interquartile range of centrality across windows~\citep{bonacich1987}.}
  {fig:centrality_boxplot}

%%------------------------------------------------------------
\subsection{Core Performance Results}

\inlinetab{Performance comparison, Universe U1 (Table~9). Primary
         economic claim: CVaR@95\% reduction vs Equal Weight, not the
         Sharpe comparison vs Standard CVaR~\citep{rockafellar2000}.
         Equal Weight raw-return advantage reflects higher volatility
         (24.4\% vs 19.1\%).}
  {tab:perf_u1}
  {\resizebox{\linewidth}{!}{%
  \begin{tabular}{lrrrrrr}
  \toprule
  Strategy & Ann.\ Ret.\ (\%) & Vol (\%) & Sharpe & MDD (\%) &
  CVaR@95\% (\%) & Sortino \\
  \midrule
  G-CVaR (ours) & 13.784 & 19.108 & 0.5644 & 73.817 & 2.849 & 0.7771 \\
  Std CVaR      & 13.764 & 19.102 & 0.5635 & 73.701 & 2.848 & 0.7758 \\
  Mean-Variance & 13.334 & 18.769 & 0.5506 & 72.196 & 2.852 & 0.7405 \\
  Equal Weight  & 16.041 & 24.404 & 0.5344 & 84.279 & 3.657 & 0.7367 \\
  \bottomrule
  \end{tabular}}}

\inlinetab{Six-strategy benchmark comparison (Table~9B): mean $\pm$ std
         across all 11 universes, 552 windows.
         HRP~\citep{lopezdeprado2016} and Risk
         Parity~\citep{maillard2010} included as concentration-aware
         baselines. G-CVaR achieves the strongest CVaR@95\% reduction
         and the highest Sharpe.}
  {tab:perf_all}
  {\resizebox{\linewidth}{!}{%
  \begin{tabular}{lrrrrrrr}
  \toprule
  Strategy & Ann.\ Ret.\ (\%) & Vol (\%) & Sharpe & Sortino & Calmar &
  MDD (\%) & CVaR@95\% (\%) \\
  \midrule
  G-CVaR (ours) & 12.818 & 16.101 & 0.6098 & 0.8364 & 0.1460 & 87.774 & 2.388 \\
  Std CVaR      & 12.869 & 15.997 & 0.6169 & 0.8504 & 0.1499 & 85.862 & 2.365 \\
  Mean-Variance & 11.462 & 15.563 & 0.5437 & 0.7185 & 0.1308 & 87.619 & 2.386 \\
  Equal Weight  & 10.706 & 20.664 & 0.3729 & 0.4922 & 0.1140 & 93.870 & 3.195 \\
  HRP           & 11.103 & 17.486 & 0.4634 & 0.6047 & 0.1212 & 91.642 & 2.713 \\
  Risk Parity   & 10.853 & 18.710 & 0.4197 & 0.5487 & 0.1172 & 92.596 & 2.903 \\
  \midrule
  \multicolumn{8}{l}{Cross-universe mean $\pm$ std:} \\
  G-CVaR   & & & $0.646 \pm 0.185$ & & & & $2.321 \pm 0.470$ \\
  Std CVaR & & & $0.653 \pm 0.184$ & & & & $2.299 \pm 0.468$ \\
  EW       & & & $0.388 \pm 0.113$ & & & & $3.132 \pm 0.515$ \\
  \bottomrule
  \end{tabular}}}

\inlinefig[0.96]{images/q1_benchmark_comparison}
  {Full six-strategy benchmark comparison across all 11 universes
   and 552 rolling windows (2005--2025). G-CVaR achieves the highest
   Sharpe ratio (0.6098) and the largest CVaR@95\% reduction (25.9\%)
   relative to Equal Weight. HRP and Risk Parity are included as
   concentration-aware baselines. Cross-universe mean $\pm$ one
   standard deviation shown~\citep{lopezdeprado2016,maillard2010}.}
  {fig:benchmark}

\textbf{Economic significance.} For a \$1,000,000 portfolio over
2005--2025, G-CVaR grows to \$11,158,395 versus \$7,645,312 for Equal
Weight — a terminal value gain of \$3,513,083. The cross-universe mean
CVaR@95\% of 2.321\% per day versus 3.132\% for Equal Weight implies a
daily tail-loss avoidance of approximately \$8,090 per \$1M
portfolio~\citep{harvey2016,bailey2012}.

\textbf{U2 Financial Services.}
The one universe where G-CVaR Sharpe (0.3435) noticeably underperforms
Std CVaR (0.3606, $\Delta = -0.0170$) is Financial Services. Large banks
(JPMorgan, Bank of America, Goldman Sachs, Wells Fargo) are core Vanguard
and BlackRock holdings across every sector, giving them consistently high
eigenvector centrality. The graph penalty therefore persistently
underweights systemically important financial stocks even during calm
periods (2012--2019) when those stocks were among the strongest performers.
This is the predictable consequence of the framework's risk-governance
mandate in a sector characterised by structurally high institutional
co-ownership~\citep{azar2018,billio2012}.

%%------------------------------------------------------------
\subsection{Crisis-Period Isolation}

\inlinetab{Crisis period isolation (Table~10, U1): MDD and CVaR@95\%
         for G-CVaR, Std CVaR, and Equal Weight during three predefined
         crisis periods~\citep{billio2012}.}
  {tab:crisis}
  {\resizebox{\linewidth}{!}{%
  \begin{tabular}{llrrr}
  \toprule
  Crisis Period & Strategy & MDD (\%) & CVaR@95\% (\%) & MDD vs EW \\
  \midrule
  2008 GFC   & G-CVaR       & 66.321 & 4.983 & $-13.0\%$ \\
             & Std CVaR     & 65.931 & 4.983 & ---        \\
             & Equal Weight & 76.162 & 6.123 & ---        \\
  \midrule
  2020 COVID & G-CVaR       & 41.419 & 5.055 & $+5.5\%$  \\
             & Std CVaR     & 41.392 & 5.056 & ---        \\
             & Equal Weight & 39.270 & 6.307 & ---        \\
  \midrule
  2022 Rate  & G-CVaR       & 26.578 & 3.391 & $-39.8\%$ \\
             & Std CVaR     & 26.484 & 3.391 & ---        \\
             & Equal Weight & 44.152 & 4.473 & ---        \\
  \bottomrule
  \end{tabular}}}

\inlinetab{Regime-separated performance (Table~13C): G-CVaR vs Equal
         Weight across all 11 universes, stratified by regime.
         \textbf{The Crisis-regime CVaR reduction of 20.7\% and MDD
         reduction of 32.5\,pp vs Equal Weight are the primary claims
         of the paper}, sourced from this table.
         $^{**}p < 0.01$; ns: not significant~\citep{harvey2016}.}
  {tab:regime_separated}
  {\resizebox{\linewidth}{!}{%
  \begin{tabular}{lrrrrrrrr}
  \toprule
  Regime & N Obs & G-CVaR Sharpe & EqW Sharpe & G-CVaR CVaR\% &
  EqW CVaR\% & CVaR Red. & G-CVaR MDD\% & EqW MDD\% \\
  \midrule
  Calm     & 112,140 & 0.9451 & 0.6985 & 1.857 & 2.508 & 26.0\% & 30.16 & 50.21 \\
  Elevated &  21,420 & 0.2174 & $-0.040$ & 3.264 & 4.348 & 24.9\% & 65.10 & 97.28 \\
  Crisis   &   5,544 & $-0.891$ & $-1.079$ & 4.423 & 5.581 & 20.7\% & \textbf{68.5} & \textbf{101.0} \\
  \midrule
  \multicolumn{3}{l}{Crisis $p$ (G-CVaR vs EqW)} & & & & $0.0000^{**}$ & & \\
  \multicolumn{3}{l}{Calm $p$ (G-CVaR vs EqW)}   & & & & $0.7677$ (ns) & & \\
  \bottomrule
  \end{tabular}}}

Table~\ref{tab:regime_separated} is the source for the abstract's 32.5\,pp
MDD figure~\citep{billio2012}: during Crisis-regime windows (5,544 daily
observations pooled across all 11 universes), G-CVaR's effective MDD of
68.5\% compares with 101.0\% for Equal Weight, a reduction of 32.5
percentage points. Calm-regime non-significance ($p = 0.7677$) is
architecturally expected — by design, $\lambda_t \approx 0$ in calm
windows~\citep{ang2002,hamilton1989}.

%%------------------------------------------------------------
\subsection{Walk-Forward Validation}
\label{sec:wfv_results}

\inlinetab{Walk-forward calibration results (Table~11, U1).
         IS Sharpe = 0.6399 is the evaluation Sharpe under frozen
         parameters ($k=10$, $I_{\text{thresh}}=0.85$), not the grid
         search best-cell value of 0.6495. OOS Sharpe = 0.8792 reflects
         the favourable post-2016 technology bull market.
         Parameters frozen post-IS following \citet{white2000}.}
  {tab:wfv}
  {\begin{tabular}{lrrrrrr}
  \toprule
  Universe & $k^*$ & $I_{\text{thresh}}^*$ (frozen) & IS Sharpe &
  OOS Sharpe & IS--OOS & Status \\
  \midrule
  U1 & 10 & 0.85 & 0.6399 & 0.8792 & $-37.4\%$ & Pass \\
  \bottomrule
  \end{tabular}}

\inlinefig[0.95]{images/plot12_walkforward_U1}
  {Walk-forward validation results for Universe U1. In-sample (IS)
   mean Sharpe = 0.6399 across 26 windows (end date $<$ 2016-01-01)
   under frozen parameters $k=10$, $I_{\text{thresh}}=0.85$.
   Out-of-sample (OOS) Sharpe = 0.8792 across 25 windows
   (start date $\geq$ 2016-01-01). Negative degradation indicates OOS
   improvement, consistent with the favourable post-2016 technology bull
   market~\citep{gu2020,harvey2016}.}
  {fig:walkforward}

The complete 416-solve grid confirms all parameter combinations produce IS
Sharpe between 0.6398 and 0.6495, a narrow range confirming robustness to
threshold choice~\citep{harvey2016}.

%%------------------------------------------------------------
\subsection{Ablation Study}

\inlinetab{Four-condition ablation (Table~12, U1). Static-$\lambda$
         ($\lambda_t = \lambda_{\max}$, real $\mathbf{c}_w$) achieves
         the highest overall Sharpe and lowest MDD. Full G-CVaR
         outperforms all partial conditions \textit{during Crisis-regime
         windows}, not necessarily on overall 51-window averages.}
  {tab:ablation}
  {\resizebox{\linewidth}{!}{%
  \begin{tabular}{lrrrrrrrr}
  \toprule
  Condition & Ann.\ Ret. & Vol & Sharpe & Sortino & Calmar &
  MDD & CVaR@95\% & Turnover \\
  \midrule
  Full G-CVaR (sigmoid + real $c_t$) & 13.784 & 19.108 & 0.5644 & 0.7771 & 0.1867 & 73.82 & 2.849 & 0.208 \\
  No Graph (sigmoid + unif.\ $c_t$)  & 13.764 & 19.102 & 0.5635 & 0.7758 & 0.1868 & 73.70 & 2.848 & 0.207 \\
  Static $\lambda$ ($\lambda_{\max}$ + real $c_t$) & 14.230 & 19.235 & 0.5838 & 0.8029 & 0.1985 & 71.70 & 2.874 & 0.191 \\
  No Both ($\lambda_t=0$, unif.\ $c_t \equiv$ Std CVaR) & 13.764 & 19.102 & 0.5635 & 0.7758 & 0.1868 & 73.70 & 2.848 & 0.207 \\
  \bottomrule
  \end{tabular}}}

\inlinefig[0.96]{images/plot13_ablation_U1}
  {Ablation study grouped bar chart (U1): Sharpe ratio and CVaR@95\%
   across the four conditions. Static-$\lambda$ achieves the highest
   overall Sharpe (0.5838) because the permanent penalty
   coincidentally tilts the portfolio away from high-volatility
   large-cap technology leaders. Full G-CVaR is predicted to
   outperform partial conditions during Crisis-regime windows
   only~\citep{ang2002,hamilton1989}.}
  {fig:ablation}

No Graph and No Both produce identical results (0.5635 Sharpe, 2.848\%
CVaR). This confirms that the centrality signal has no effect when
$\lambda_t \approx 0$, which occurs in 38 of 51 windows — the
regime-sparsity property operating exactly as intended~\citep{ang2002}.

%%------------------------------------------------------------
\subsection{Statistical Significance and Effect Size}

\inlinetab{Pooled cross-universe statistical tests (Table~13B): 11
         universes, 139,104 paired daily observations. Wilcoxon
         signed-rank used as primary test (Shapiro-Wilk rejects
         normality for all universes~\citep{cont2001}).
         $^{**}p < 0.01$; ns: not significant at $\alpha = 0.05$.}
  {tab:stats}
  {\resizebox{\linewidth}{!}{%
  \begin{tabular}{llrrrrl}
  \toprule
  Comparison & Metric & Wilcoxon $p$ & $t$-test $p$ & Cohen's $d$ &
  Bootstrap 95\% CI & Sig. \\
  \midrule
  G-CVaR vs Std CVaR   & $\Delta$Sharpe & 0.0517 & 0.4328 & $-0.0021$ &
  $[-0.0148, +0.0011]$ & ns \\
  G-CVaR vs Equal Wt   & $\Delta$Sharpe & 0.0065 & 0.0000 & $+0.0143$ &
  $[+0.1998, +0.2739]$ & $^{**}$ \\
  \bottomrule
  \end{tabular}}}

The Wilcoxon $p = 0.0517$ for G-CVaR versus Standard CVaR does not
achieve the conventional 0.05 threshold and must be reported precisely —
no rounding is appropriate~\citep{harvey2016}. This non-significance is
architecturally expected: in 38 of 51 U1 windows (74.5\%), $\lambda_t
\approx 0$ by design~\citep{ang2002,hamilton1989}.

\inlinefig[0.96]{images/q1_statistical_validation}
  {Q1 statistical validation package (11 universes). \textit{Left:}
   bootstrap distribution of $\Delta$Sharpe (G-CVaR vs Std CVaR)
   centred near zero, consistent with regime-sparsity architecture.
   \textit{Centre:} CVaR@95\% by regime — G-CVaR outperforms Equal
   Weight in all three regimes; the Crisis-regime gap is most
   pronounced ($p = 0.0000$). \textit{Right:} effect size per
   universe (Cohen's $d$), showing the range $[-0.0114, +0.0043]$
   across 11 universes~\citep{harvey2016,cont2001}.}
  {fig:stat_validation}

\inlinefig[0.82]{images/plot14_kde_U1}
  {Kernel density comparison of daily returns: G-CVaR vs Standard
   CVaR (Universe U1). The two distributions are nearly
   indistinguishable in calm-regime windows, confirming that the
   graph penalty contributes only when $\lambda_t > 0$~\citep{cont2001}.}
  {fig:return_kde}

\inlinetab{Universe-level Wilcoxon $p$-values: G-CVaR vs Std CVaR
         $\Delta$Sharpe (Table~13A). Zero of 11 universes achieve
         $p < 0.05$. Architecturally expected and disclosed
         explicitly~\citep{ang2002,harvey2016}.}
  {tab:stats_universe}
  {\resizebox{\linewidth}{!}{%
  \begin{tabular}{lrrrrrr}
  \toprule
  Universe & N & G-CVaR Sharpe & StdCVaR Sharpe & $\Delta$Sharpe &
  Wilcoxon $p$ & Cohen's $d$ \\
  \midrule
  U1  & 12,852 & 0.5644 & 0.5635 & $+0.0009$ & 0.2537 & $+0.0040$ \\
  U2  & 12,852 & 0.3435 & 0.3606 & $-0.0170$ & 0.5414 & $-0.0114$ \\
  U3  & 12,600 & 0.4532 & 0.4690 & $-0.0158$ & 0.1189 & $-0.0091$ \\
  U4  & 12,348 & 0.7779 & 0.7905 & $-0.0126$ & 0.5937 & $+0.0032$ \\
  U5  & 12,852 & 0.8869 & 0.8962 & $-0.0093$ & 0.8502 & $-0.0031$ \\
  U6  & 12,852 & 0.6075 & 0.6214 & $-0.0139$ & 0.4639 & $-0.0010$ \\
  U7  & 12,348 & 1.0122 & 1.0207 & $-0.0085$ & 0.0628 & $-0.0024$ \\
  U8  & 12,348 & 0.5402 & 0.5400 & $+0.0002$ & 0.8565 & $+0.0009$ \\
  U9  & 12,852 & 0.7421 & 0.7450 & $-0.0029$ & 0.4142 & $-0.0038$ \\
  U10 & 12,348 & 0.6295 & 0.6254 & $+0.0041$ & 0.9045 & $+0.0043$ \\
  U11 & 12,852 & 0.5447 & 0.5529 & $-0.0081$ & 0.2696 & $-0.0035$ \\
  \bottomrule
  \end{tabular}}}

%%------------------------------------------------------------
\subsection{Monte Carlo Stress Testing}

\inlinetab{Monte Carlo survival statistics (Table~14, U1, 10,500 paths).
         100\% survival is expected for any diversified portfolio under
         single-period Gaussian scenarios~\citep{jagannathan2003,glasserman2004};
         the meaningful comparison is CVaR@95\% across strategies.}
  {tab:montecarlo}
  {\begin{tabular}{lrrrr}
  \toprule
  Strategy & Mean MDD & P95 MDD & CVaR@95\% & Survival $<30\%$ \\
  \midrule
  G-CVaR       & 0.96\% & 3.88\% & 5.42\% & 100\% \\
  Std CVaR     & 0.96\% & 3.86\% & 5.40\% & 100\% \\
  Mean-Variance& 0.95\% & 3.83\% & 5.37\% & 100\% \\
  Equal Weight & 1.12\% & 4.57\% & 6.36\% & 100\% \\
  \bottomrule
  \end{tabular}}

\inlinefig[0.92]{images/plot15_montecarlo_U1}
  {Monte Carlo CVaR@95\% comparison across 10,500 GFC-calibrated
   synthetic paths (U1). G-CVaR (5.42\%) achieves a 14.8\%
   tail-loss reduction versus Equal Weight (6.36\%), validating
   the GFC-covariance-based stress test as a reasonable proxy for
   crisis-window behaviour~\citep{glasserman2004,cont2001}.}
  {fig:montecarlo}

%%------------------------------------------------------------
\subsection{Sensitivity Analysis}

\inlinefig[0.88]{images/plot16_ithresh_heatmap_U1}
  {Sharpe ratio sensitivity to $I_{\text{thresh}}$ ($k=10$ fixed):
   values 0.75--0.95 produce Sharpe in the range 0.7567--0.7608,
   a variation of only 0.0041 Sharpe points. Confirms robustness
   to precise threshold choice~\citep{kolm2014,harvey2016}.}
  {fig:sensitivity_ithresh}

\inlinefig[0.88]{images/plot17_k_heatmap_U1}
  {Sharpe ratio sensitivity to sigmoid steepness $k$ ($I_{\text{thresh}}
   = 0.85$ fixed): values 5--20 produce Sharpe 0.7578--0.7635.
   Narrow range confirms robustness to $k$ choice. The frozen
   walk-forward values ($k=10$, $I_{\text{thresh}}=0.85$) are used
   throughout all reported results~\citep{white2000,harvey2016}.}
  {fig:sensitivity_k}

%%------------------------------------------------------------
\subsection{Regime Sparsity Disclosure}

For U1, only 2 of 51 windows (3.9\%) reach Crisis regime. The mean
$\lambda_t$ in calm windows is 0.00502, rising to 0.81266 in crisis
windows. Across all 11 universes, 22 of 552 total windows (4.0\%) reach
Crisis. This sparsity is intentional~\citep{ang2002,hamilton1989}:
penalising contagion risk in calm markets would impose permanent return
drag without commensurate crisis-protection benefit.

\inlinefig[0.95]{images/plot18_activation_heatmap}
  {$\lambda_t$ activation heatmap across all 11 universes and 552
   windows. Crisis-regime windows (dark cells, $\lambda_t > 0.5$)
   account for only 4.0\% of total windows. Universes with fewer
   than 5 crisis activations are flagged per Q1 reporting
   standards~\citep{harvey2016}.}
  {fig:activation_heatmap}

%%------------------------------------------------------------
\subsection{XAI Attribution Fidelity}

\inlinetab{Top-10 most penalised stocks (Table~17, U1): mean eigenvector
         centrality, mean weight delta, and appearance count.
         Negative $\Delta w$ confirms penalty operates in the predicted
         direction~\citep{bonacich1987,arrieta2020}.}
  {tab:top10_penalised}
  {\begin{tabular}{llrrr}
  \toprule
  Rank & Ticker & Mean $c_t$ & Mean $\Delta w$ & Appearances \\
  \midrule
  1  & MSFT & 0.24675 & $-0.2309$ & 16 \\
  2  & QCOM & 0.24191 & $-0.1608$ & 24 \\
  3  & ADI  & 0.25219 & $-0.1245$ & 33 \\
  4  & INTU & 0.25529 & $-0.0714$ & 23 \\
  5  & CRM  & 0.24434 & $-0.0193$ & 19 \\
  6  & CSCO & 0.24486 & $-0.0076$ & 16 \\
  7  & AMD  & 0.24121 & $+0.0000$ & 3  \\
  8  & LRCX & 0.25768 & $+0.0000$ & 3  \\
  9  & NVDA & 0.24555 & $+0.0000$ & 4  \\
  10 & KLAC & 0.24530 & $+0.0000$ & 8  \\
  \bottomrule
  \end{tabular}}

\inlinefig[0.92]{images/plot19_xai_attribution_U1}
  {XAI attribution scatter: eigenvector centrality $c_t$ vs weight
   delta $\Delta w = w^{\text{G-CVaR}} - w^{\text{Std-CVaR}}$ for
   all 51 windows (U1). Pearson $r = -0.510$ during Elevated-regime
   window W08 ($p = 0.036$) is the strongest observed correlation,
   confirming attribution fidelity is highest when the penalty is
   most active~\citep{bonacich1987,arrieta2020}.}
  {fig:xai_attribution}

%%------------------------------------------------------------
\subsection{HITL Governance Decisions}

\inlinetab{HITL decision summary (Table~21, U1). The Constrain action
         re-runs G-CVaR with $w_{\max} = 8\%$, enforcing tighter
         position limits.}
  {tab:hitl_summary}
  {\begin{tabular}{lrrrr}
  \toprule
  Universe & Total Events & Approved & Rejected & Constrained \\
  \midrule
  U1 & 11 & 5 & 3 & 3 \\
  \bottomrule
  \end{tabular}}

\inlinefig[0.96]{images/plot21_governance_report_U1}
  {Sample governance report from the Gradio HITL interface for
   Window W09 ($I_t = 1.000$, Crisis trigger). The report presents
   the trigger reason, proposed G-CVaR and Std-CVaR weights,
   weight deltas, centrality ranks, and XAI narrative before any
   action is taken~\citep{holzinger2016,cai2019}.}
  {fig:governance_report}

%%------------------------------------------------------------
\subsection{System Trustworthiness Evaluation}

\inlinetab{System trustworthiness evaluation (Table~27): four criteria
         across 160 governance scenarios from all 11 universes.
         Zero-hallucination reported with both strict and substantive
         interpretations~\citep{ji2023,arrieta2020}.}
  {tab:trust}
  {\begin{tabular}{lrrrr}
  \toprule
  Criterion & Score & Threshold & Pass \\
  \midrule
  Trigger Accuracy              & 100.0\% (160/160) & $\geq 90\%$ & \checkmark \\
  Weight Conservatism           & 100.0\% (52/52)   & $\geq 80\%$ & \checkmark \\
  Narrative Accuracy            &  90.6\% (145/160) & $\geq 70\%$ & \checkmark \\
  Zero-Hallucination (strict)   &  70.0\% (35/50)   & $\geq 70\%$ & \checkmark \\
  Zero-Hallucination (subst.)   & 100.0\% (50/50)   & N/A          & \checkmark \\
  \bottomrule
  \end{tabular}}

\inlinefig[0.82]{images/plot20_radar_U1}
  {Agent~4 trustworthiness evaluation radar chart (Table~27, U1,
   160 governance scenarios). All four criteria pass their respective
   thresholds. Narrative Accuracy (90.6\%) is the primary XAI
   faithfulness metric; zero-hallucination (70.0\% strict /
   100.0\% substantive) reflects formatting-precision mismatches,
   not factual financial errors~\citep{arrieta2020,doshi2017,ji2023}.}
  {fig:trustworthiness_radar}

%%------------------------------------------------------------
\subsection{HITL Governance Ablation}

\inlinetab{HITL governance ablation (Table~26): four conditions across
         all 11 universes, 552 windows, 139,104 daily observations.
         Cost of governance (0.1264 Sharpe, 3.12\,pp MDD) represents
         the price of regulatory-compliant oversight.}
  {tab:hitl_ablation}
  {\resizebox{\linewidth}{!}{%
  \begin{tabular}{lrrrr}
  \toprule
  Condition & Sharpe & MDD (\%) & CVaR@95\% (\%) & Ann.\ Ret.\ (\%) \\
  \midrule
  Full Governed (G-CVaR + HITL Constrain) & 0.4834 & 90.894 & 2.664 & 11.458 \\
  Auto G-CVaR (no HITL)                  & 0.6098 & 87.774 & 2.388 & 12.818 \\
  Std CVaR (no governance)               & 0.6169 & 85.862 & 2.365 & 12.869 \\
  Equal Weight (no oversight)            & 0.3729 & 93.870 & 3.195 & 10.706 \\
  \bottomrule
  \end{tabular}}}

\inlinefig[0.96]{images/plot25_hitl_ablation}
  {HITL governance ablation cumulative return curves (11 universes,
   552 windows). Full Governed MDD (90.894\%) exceeds Auto G-CVaR
   (87.774\%) by 3.12\,pp: the Constrain action redistributes
   weight away from high-return concentrated positions during bull
   markets, quantifying the cost of regulatory compliance~\citep{nystrup2019}.}
  {fig:hitl_ablation}

%%------------------------------------------------------------
\subsection{MAS Architecture Validation}

\inlinetab{MAS vs monolithic architecture (Table~28, U1, 10 windows,
         5 runs). The 1.267-second overhead per 10-window batch
         purchases complete audit trail and fault isolation~\citep{hayesroth1985,wooldridge2009}.}
  {tab:mas}
  {\resizebox{\linewidth}{!}{%
  \begin{tabular}{lrrrrrr}
  \toprule
  Architecture & Mean Lat.\ (s) & Lat.\ Std & Error Recovery &
  Fault Logged & Stale Det. & Audit Trail \\
  \midrule
  MAS (Blackboard)   & 3.886 & 2.495 & 100\% & \checkmark & \checkmark & Full \\
  Monolithic         & 2.619 & 0.198 & 100\% & $\times$   & $\times$   & None \\
  Direct Agent Calls & 2.545 & 0.151 & 100\% & $\times$   & $\times$   & None \\
  \bottomrule
  \end{tabular}}}

\inlinefig[0.90]{images/plot_latency_U1}
  {Per-agent empirical latency breakdown (Table~23, U1, averaged
   over 5 windows). Agent~3 (G-CVaR + Std CVaR solve) dominates
   at 59.0\% of total pipeline time (0.1190\,s/window), consistent
   with $\mathcal{O}(N^2 T)$ theoretical complexity~\citep{rockafellar2000,diamond2016}.}
  {fig:latency}

\inlinefig[0.96]{images/plot26_mas_validation}
  {MAS architecture validation: latency distribution across 5 runs
   for MAS (blackboard), monolithic pipeline, and direct agent calls
   (U1, 10-window batch). MAS overhead of 1.267\,s purchases fault
   isolation, stale-data detection, and full audit trail — structural
   advantages unavailable in monolithic designs at any
   latency~\citep{hayesroth1985,erman1980,mifid2014}.}
  {fig:mas_validation}

%%------------------------------------------------------------
\subsection{Cross-Universe Robustness}

\inlinetab{Cross-universe performance (Table~9A): G-CVaR, Standard CVaR,
         and Equal Weight across all 11 universes. G-CVaR outperforms
         Equal Weight in Sharpe in all 11 universes; U7 Industrials
         and U4 Energy show the largest advantages~\citep{harvey2016}.}
  {tab:cross_universe}
  {\resizebox{\linewidth}{!}{%
  \begin{tabular}{llrrrrr}
  \toprule
  Universe & Sector & G-CVaR Sharpe & StdCVaR Sharpe & EqW Sharpe &
  G-CVaR CVaR\% & EqW CVaR\% \\
  \midrule
  U1  & Technology         & 0.5644 & 0.5635 & 0.5344 & 2.849 & 3.657 \\
  U2  & Financial Services & 0.3435 & 0.3606 & 0.2421 & 3.304 & 4.058 \\
  U3  & Healthcare         & 0.4532 & 0.4690 & 0.3568 & 2.167 & 2.906 \\
  U4  & Energy             & 0.7779 & 0.7905 & 0.2759 & 2.357 & 3.514 \\
  U5  & Consumer Staples   & 0.8869 & 0.8962 & 0.4950 & 1.626 & 2.336 \\
  U6  & Consumer Disc.     & 0.6075 & 0.6214 & 0.2423 & 2.449 & 3.536 \\
  U7  & Industrials        & 1.0122 & 1.0207 & 0.5167 & 1.628 & 2.452 \\
  U8  & Materials          & 0.5402 & 0.5400 & 0.2873 & 2.053 & 3.035 \\
  U9  & Utilities          & 0.7421 & 0.7450 & 0.4229 & 2.110 & 2.765 \\
  U10 & Real Estate        & 0.6295 & 0.6254 & 0.3522 & 2.551 & 3.401 \\
  U11 & Comm.\ Services   & 0.5447 & 0.5529 & 0.5374 & 2.440 & 2.793 \\
  \midrule
  \multicolumn{2}{l}{Mean $\pm$ Std} &
  $0.646 \pm 0.185$ & $0.653 \pm 0.184$ & $0.388 \pm 0.113$ &
  $2.321 \pm 0.470$ & $3.132 \pm 0.515$ \\
  \bottomrule
  \end{tabular}}}

\inlinefig[0.96]{images/plot_sector_weights}
  {Cross-sector allocation weights from the two-level optimisation
   (Phase~16): Level-2 mean-variance allocation across 11 sector
   universes over 49 windows with full coverage. Sector weights
   reflect relative Sharpe-to-volatility ratios across
   universes~\citep{markowitz1952,kolm2014}.}
  {fig:sector_weights}
